\section{Convolutional Neural Networks (CNNs)}

The Python part of this exercise is available in \myhref{https://github.com/LosDelDGIIM/LosDelDGIIM.github.io/blob/main/subjects/Erasmus-DUE/GKI/Exercises/Exercise10.ipynb}{this Jupyter Notebook}.\\

\begin{ejercicio}[Convolutional Layer]
    What is the main task of a convolutional layer and how does it help the network find patterns (like edges or shapes) in an image?

    The main task of a convolutional layer is to extract features from the input data (such as images) by applying convolution operations using filters (kernels). These filters slide over the input image and perform element-wise multiplications and summations, allowing the network to detect local patterns such as edges, textures, and shapes. By learning different filters during training, the convolutional layer can capture various features at different levels of abstraction, which helps the network understand and recognize complex patterns in the data.
\end{ejercicio}

\begin{ejercicio}[Filter]
    What does a filter (kernel) do in a convolutional layer as it moves step-by-step across the image?

    A filter (kernel) in a convolutional layer is a small matrix of weights that is applied to the input image through a process called convolution. As the filter moves step-by-step (or slides) across the image, it performs element-wise multiplications between the filter values and the corresponding pixel values of the image, followed by summing these products to produce a single output value. This operation allows the filter to detect specific features in different regions of the image, such as edges, corners, or textures, depending on the learned weights of the filter.
\end{ejercicio}

\begin{ejercicio}[Image Channels]
    What is the difference between a grayscale image and a color image (RGB) regarding the number of channels the network has to process?

    A grayscale image has a single channel, where each pixel represents the intensity of light (brightness) at that point, typically ranging from black to white. In contrast, a color image (RGB) has three channels corresponding to the red, green, and blue components of each pixel. Each channel contains intensity values for its respective color, allowing the network to process and learn from the combined information of all three channels to capture color features in addition to spatial patterns.
\end{ejercicio}

\begin{ejercicio}[Stride]
    What does the term stride mean and how does the output size change when you increase the stride?

    The term stride refers to the number of pixels by which the filter moves (or slides) across the input image during the convolution operation. A stride of 1 means the filter moves one pixel at a time, while a stride of 2 means it moves two pixels at a time. Increasing the stride reduces the spatial dimensions of the output feature map because fewer positions are sampled from the input image. As a result, larger strides lead to smaller output sizes, which can help reduce computational complexity but may also result in loss of spatial information.
\end{ejercicio}

\begin{ejercicio}[Padding]
    Why do we often add an artificial border (e.g., made of zeros) around the image before applying convolution?

    We often add an artificial border (e.g., made of zeros) around the image given two main reasons:
    \begin{itemize}
        \item To preserve the spatial dimensions of the output feature map. Without padding, the size of the feature map decreases after each convolution operation, which can lead to a loss of information, especially in deeper networks.
        \item To ensure that the filter can be applied to the edges of the image. Without padding, the filter has fewer neighboring pixels to convolve with at the borders, which can lead to artifacts in the feature map and may not capture important edge features.
    \end{itemize}
    Padding helps maintain the integrity of the input data and allows the network to learn features from the entire image, including the edges.
\end{ejercicio}

\begin{ejercicio}[Max Pooling]
    How does max pooling work and why does it help reduce the amount of data?

    Max pooling is a downsampling operation that reduces the spatial dimensions of the feature maps while retaining the most important information. It works by dividing the input feature map into non-overlapping regions (usually squares) and selecting the maximum value from each region to create a smaller output feature map. This process helps reduce the amount of data by summarizing the most prominent features in each region, which can lead to lower computational complexity, reduced memory usage, and improved robustness to variations in the input data, such as translations and distortions.
\end{ejercicio}

\begin{ejercicio}[Purpose of Pooling]
    Why does a pooling layer almost always follow directly after a convolutional layer in a CNN architecture?

    A pooling layer almost always follows directly after a convolutional layer in a CNN architecture because it serves to downsample the feature maps produced by the convolutional layer. This downsampling reduces the spatial dimensions of the feature maps, which helps to decrease computational complexity and memory usage. Additionally, pooling layers help to make the network more robust to variations in the input data, such as translations and distortions, by retaining the most important features while discarding less relevant information. By following a convolutional layer with a pooling layer, the network can effectively capture hierarchical features at different levels of abstraction.
    % // TODO: Why should we add pooling when we can increase the stride?
\end{ejercicio}

\begin{ejercicio}[Flattening]
    What happens during flattening and why must the data be made flat before it flows into the final layers?

    During flattening, the multi-dimensional feature maps produced by the convolutional and pooling layers are transformed into a one-dimensional vector. This process is necessary because the final layers of a CNN, typically fully connected (dense) layers, expect input in a flat format. Flattening allows the network to combine the learned features from the previous layers and pass them to the fully connected layers for further processing, such as classification or regression. By converting the spatially structured data into a flat vector, the network can effectively learn complex relationships between the features and make predictions based on the extracted information.
\end{ejercicio}

\begin{ejercicio}[Fully Connected Layer]
    What is the role of the classic, fully connected layers (dense layers) at the very end of a CNN?

    The role of the classic, fully connected layers (dense layers) at the very end of a CNN is to perform high-level reasoning and make final predictions based on the features extracted by the preceding convolutional and pooling layers. These layers take the flattened feature vector as input and learn complex relationships between the features through weighted connections. The fully connected layers combine the learned features to produce output values, such as class probabilities in classification tasks or continuous values in regression tasks. They serve as the decision-making part of the network, translating the learned representations into actionable outputs.
\end{ejercicio}

\begin{ejercicio}[Number of Filters]
    Why do we usually use more filters in the deeper layers of a CNN compared to the very first layers?

    We usually use more filters in the deeper layers of a CNN compared to the very first layers because deeper layers are responsible for capturing more complex and abstract features of the input data. The initial layers typically learn simple features such as edges and textures, which can be represented with fewer filters. As we move deeper into the network, the features become more intricate, requiring a larger number of filters to capture various combinations of lower-level features and to represent higher-level concepts. By increasing the number of filters in deeper layers, the network can learn a richer set of features, improving its ability to recognize complex patterns and make accurate predictions.
\end{ejercicio}